\section{Positive-source compression and prior spectral results}\label{vfourD:section}
This appendix incorporates the source-bank argument of \cite{vfourSourceCompression} as an alternative proof for the same qualitative spectral conclusions previously obtained in \cite{vfourSchubr}. The earlier result and its relationship to the current realization are stated in Section~\ref{vfourD:prior}. No historical priority or better-constant claim is made. Unlike the convention used in the main text, inner products in this appendix are conjugate-linear in the first argument.

We retain the actual form $q_0$, $\Hh$, $X$, $m$ and weight $W$ from Section~\ref{a:sec:energy}. Let $\Hh_c=\ker m$, $X_c=X\cap\Hh_c$, $C u=u-2m(u)\kappa$, and let $L_c$ be the form representative of $q_0$ on $X_c$. All identities used below are on this complete centered form domain. The full theta normalization and envelope give
\begin{equation}\label{vfourD:Wbound}
 \norm\kappa_{\Hh}=1,\qquad B_0:=\int_\R W(x)\,dx\le\tfrac14,\qquad
 0<W(x)\le40e^{4|x|-\pi e^{2|x|}}\le40e^{4|x|-3e^{2|x|}}.
\end{equation}
The Gamma symbol $d_\Gamma$ and $c_\Gamma$ are those of the complete form; $d_\Gamma(t)\ge\frac12\log(1+4t^2)$, and every prime power and both shift orientations are retained.

For integer $N\ge1$, write
\begin{gather*}
 c_N(t)=c_\Gamma+2\sum_{2\le n\le N}\frac{\Lambda(n)}{\sqrt n}\cos(t\log n),\quad
 C_N=c_\Gamma+2\sum_{2\le n\le N}\frac{\Lambda(n)}{\sqrt n},\\
 a_N=(d_\Gamma-c_N)_+,\quad b_N=(c_N-d_\Gamma)_+,\quad
 T_N=\tfrac12\sqrt{e^{2C_N}-1}>0.
\end{gather*}
The scalar $b_N$ is nonnegative, even, vanishes for $|t|\ge T_N$, and satisfies
\begin{equation}\label{vfourD:b-integral}
 \int_0^{T_N}b_N(t)dt\le T_N-\tfrac12\arctan(2T_N)\le T_N.
\end{equation}
For $g_t(x)=W(x)\cos(tx)/\sqrt{2\pi}$, the actual positive trace-class deficit on $\Hh_c$ is
\begin{equation}\label{vfourD:Bform}
 B_N=2\int_0^{T_N}b_N(t)(Cg_t)(Cg_t)^*dt,\qquad
 q_0[u]=\int a_N(t)|\widehat u(t)|^2dt-\ip u{B_Nu}-\ip u{D_{>N}u}
\end{equation}
for every $u\in X_c$. The complete centered tail is bounded self-adjoint and
\begin{equation}\label{vfourD:tail}
 \norm{D_{>N}}\le\tau_N\le160(N+1)^2e^{-\pi(N+1)}.
\end{equation}
The full-domain form identity follows from finite-head Plancherel and the entire prime-tail norm convergence; the last two terms retain their actual signs. In particular,
\begin{equation}\label{vfourD:lower}
 q_0[u]\ge-\ip u{B_Nu}-\tau_N\norm u_{\Hh}^2\qquad(u\in X_c).
\end{equation}
For $\delta>0$ define the \emph{closed} negative count
\[N_-(\delta)=\dim\operatorname{Ran}\mathbf1_{(-\infty,-\delta]}(L_c),\]
counting eigenvalues at $-\delta$ and multiplicities. A publicly posted preceding trace paper \cite{vfourTrace}, which uses the same actual input realization and classical trace-count principle, proves $N_-(\delta)\le\delta^{-1}e^{O(\sqrt{\log(1/\delta)})}$ and $(L_c)_-\in\mathfrak S_p$ for $p>1$. Earlier, Schubr \cite[Theorem~1.1]{vfourSchubr} obtained a subpolynomial negative-count estimate and all positive Schatten exponents for a constrained weighted Weil operator. Section~\ref{vfourD:prior} explains the even-sector form identification. The argument below is an alternative proof with an explicit constant, not a claim of first discovery or of a better constant.

\subsection{Relation to earlier work}\label{vfourD:prior}
Schubr's publicly posted edition~3 \cite{vfourSchubr}, Theorem~1.1, gives a subpolynomial negative spectral count and membership of the negative part in every $\mathfrak S_p$, $p>0$, for the complete constrained operator $H$ defined there. We record the connection with the present form, rather than claiming those consequences as first results here.

In \cite[Section~2, equations (4)--(6)]{vfourSchubr}, the source $\psi$ is our $\kappa$, the weight $w=\psi\cosh(x/2)$, $s^2=\psi/\cosh(x/2)=2W$, the multiplier $a$ equals $d_\Gamma$, and the scalar $m_0$ there is $-c_\Gamma$. The unitary map $U_W:\Hh\to L^2_{\rm even}(dx)$ is $U_Wu=u/\sqrt W$. For even $u$, the second constraint of \cite{vfourSchubr} is automatic by parity, while its first is $\langle\sqrt w,U_Wu\rangle=\sqrt2\,m(u)$. Thus $U_W\Hh_c$ is its even constrained space. On centered even compact-smooth functions, the full form of \cite{vfourSchubr} at $U_Wu$ is exactly $2q_0[u]$: $sU_Wu=\sqrt2u$, both complete prime orientations are retained, and its nonunitary Fourier transform with $dt/(2\pi)$ matches our unitary convention.

The form closures agree under the stated source inputs. Multiplication by $\sqrt W$ or its reciprocal maps compact smooth even functions bijectively to compact smooth even functions because $W$ is smooth and strictly positive. The ambient compact-smooth form core in Lemma~\ref{a:lem:Xcore} yields a centered compact-smooth core: after approximating $u\in X_c$, subtract from each approximant a scalar multiple of one fixed even compact-smooth test with nonzero $m$; the scalar tends to zero by continuity of $m$. This is precisely the transformed even constrained core in \cite{vfourSchubr}. Consequently the even restriction of its form-defined operator is unitarily equivalent to $2L_c$. The full constrained operator there reduces under parity, so its negative-count and all-Schatten estimates imply the corresponding conclusions for $L_c$; the use of strict thresholds there versus the closed threshold here does not alter the asymptotic, since the closed $L_c$ threshold $-\delta$ is contained in its strict range below $-\delta/2$, corresponding to the strict $H$ threshold $-\delta$. The explicit numerical constant $40$ below is one admissible bound, not a claim to improve the constants of \cite{vfourSchubr}.

\subsection{A centered source bank and its true Gram}\label{vfourD:bank}
Fix $R>0$ and integer $m\ge0$. In $\Hh_c$ set
\begin{equation}\label{vfourD:bankdef}
 v_j=\frac1{\sqrt{2\pi}}C\!\left(W(x)x^{2j}\mathbf1_{[-R,R]}(x)\right),\quad 0\le j\le m,\quad V=(v_0,\ldots,v_m),\quad r=m+1.
\end{equation}
Each vector belongs to $\Hh_c$, and the bank is independent. Indeed, if $C(Wp\mathbf1_{[-R,R]})=0$ for an even polynomial $p$, the uncentered vector is a scalar multiple of $\kappa$. Outside the interval it vanishes while $\kappa>0$, so the scalar is zero; then $W>0$ on the interval forces $p=0$. This proof applies to complex coefficients.

Let
\begin{equation}\label{vfourD:gram}
 a_j=\int_{-R}^{R}\kappa(x)x^{2j}dx,\qquad w_\ell=\int_{-R}^{R}W(x)x^{2\ell}dx,\qquad
 G_{ij}=\ip{v_i}{v_j}=\frac{w_{i+j}-a_i a_j}{2\pi}.
\end{equation}
Thus $G=V^*V$ is positive definite as a matrix. The integrals in \eqref{vfourD:gram} are over the indicated \emph{finite} interval, and the centering terms are retained. Define
\[P=VG^{-1}V^*,\qquad Q=I-P\quad\text{on }\Hh_c.\]
This is the true weighted orthogonal projection, not a coefficient-identity Gram or a projection onto the physical kernel. The source bank is required only to lie in $\Hh_c$. We do not require or use $v_j\in X_c$, $v_j\in\dom L_c$, or preservation of either domain by $P$ or $Q$.

\subsection{Trace defect and an explicit rank}\label{vfourD:compression}
\begin{theorem}\label{vfourD:defect}
For every $\varepsilon>0$, with $T=T_N$, define
\begin{equation}\label{vfourD:parameters}
 h=\max\{1,T/\varepsilon\},\quad R=\tfrac12\log\log(80h),\quad
 r=\left\lceil\max\left\{1,3TR,\tfrac12\log_2(2h)\right\}\right\rceil,\quad m=r-1.
\end{equation}
For the bank \eqref{vfourD:bankdef} and $Q$ above,
\begin{equation}\label{vfourD:traceerror}
 \operatorname{tr}(QB_NQ)<\varepsilon,\qquad
 r\le3TR+\tfrac12\log_2(2h)+2.
\end{equation}
\end{theorem}
\begin{proof}
Set $k=2m+2=2r$ and $p_{m,t}(x)=\sum_{j=0}^m(-1)^j(tx)^{2j}/(2j)!$. For $|x|\le R$ and $0\le t\le T$, Taylor's theorem to order $2m+1$ gives
\[|\cos(tx)-p_{m,t}(x)|\le E:=\frac{(TR)^k}{k!}.
\]
The centered truncated Taylor vector belongs to $\operatorname{Ran}P$. Since $C$ is contractive, the projection distance, with disjoint inside and outside regions, satisfies
\begin{equation}\label{vfourD:distance}
 \norm{QCg_t}_{\Hh}^2\le\frac{W_{\rm tail}(R)+B_0E^2}{2\pi},\qquad W_{\rm tail}(R)=\int_{|x|>R}W(x)dx.
\end{equation}
Consequently, by \eqref{vfourD:Bform} and \eqref{vfourD:b-integral},
\begin{equation}\label{vfourD:trace-majorant}
 \operatorname{tr}(QB_NQ)=2\int_0^T b_N(t)\norm{QCg_t}_{\Hh}^2dt
 \le\frac{T}{\pi}\left(W_{\rm tail}(R)+B_0E^2\right).
\end{equation}
This bound includes the full centered tail of each source: no uncentered truncation is treated as already centered.

For $Y=e^{2R}\ge1$, the envelope in \eqref{vfourD:Wbound} has logarithmic derivative at most $-2$ on $x\ge0$. Integrating the tangent exponential majorant on both half-lines yields
\begin{equation}\label{vfourD:Wtail}
 W_{\rm tail}(R)\le40Y^2e^{-3Y}\le40e^{-Y}.
\end{equation}
With \eqref{vfourD:parameters}, $Y=\log(80h)$, so $W_{\rm tail}(R)\le1/(2h)$. The factorial estimate $k!\ge(k/e)^k$, $e<3$, and $k=2r\ge6TR$ give $E\le2^{-k}\le1/(2h)$. Using $B_0\le1/4$ in \eqref{vfourD:trace-majorant},
\[\operatorname{tr}(QB_NQ)\le\frac{T}{\pi}\left(\frac1{2h}+\frac1{16h^2}\right)
 \le\frac{9T}{16\pi h}\le\frac{9\varepsilon}{16\pi}<\varepsilon.
\]
The rank bound follows from $\lceil\max(a,b,c)\rceil\le a+b+c+2$ for the displayed nonnegative quantities.
\end{proof}
For a certified finite implementation, use rigorous upper bounds for $T,h,R$ and choose an even $k$ satisfying the displayed inequalities strictly; if the radius is enlarged, recheck $k\ge6TR$. This avoids deciding whether a computable real is exactly an integer. It does not provide a bound on the condition number of $G$, required arithmetic precision, or total runtime.

\subsection{Negative count and Schatten classes}\label{vfourD:count}
\begin{theorem}\label{vfourD:main}
For $0<\delta\le1$, the actual centered negative spectral count satisfies
\begin{equation}\label{vfourD:countbound}
 N_-(\delta)<\exp\!\left(40\sqrt{1+\log(1/\delta)}\right).
\end{equation}
In particular $N_-(\delta)\le\exp(O(\sqrt{1+\log(1/\delta)}))$ as $\delta\downarrow0$. The constant $40$ is deliberately loose.
\end{theorem}
\begin{proof}
Put
\begin{equation}\label{vfourD:Ndelta}
 M=\max\{8,\lceil\log_2(2/\delta)\rceil\},\qquad N=M-1.
\end{equation}
Then $\tau_N\le\delta/2$. Indeed, for $M\ge8$,
\[160M^2e^{-\pi M}\le e^{-\pi M/2}\le2^{-M}\le\delta/2.
\]
The first inequality holds at $M=8$ because $160\cdot8^2<(5/2)^{12}<e^{3\cdot8/2}<e^{\pi8/2}$, and its ratio improves for $M\ge8$; the second follows from $\pi/2>\log2$.

Use Theorem~\ref{vfourD:defect} for $B_N$ with $\varepsilon=\delta/4$. If the closed negative spectral subspace had dimension greater than $r$, it would contain a nonzero $u\in\operatorname{Ran}Q$. Lower boundedness of $L_c$ places that spectral subspace in $\dom L_c\subset X_c$. Since $B_N\ge0$,
\[\ip u{B_Nu}=\ip u{QB_NQu}\le\operatorname{tr}(QB_NQ)\norm u_{\Hh}^2<\frac\delta4\norm u_{\Hh}^2.
\]
Then \eqref{vfourD:lower} gives $q_0[u]>-3\delta\norm u_{\Hh}^2/4$, contradicting $q_0[u]\le-\delta\norm u_{\Hh}^2$. Thus $N_-(\delta)\le r$. This argument uses only that $P$ is a bounded finite-rank projection on $\Hh_c$; it never applies $L_c$ to a bank vector or assumes $P X_c\subset X_c$.

For the displayed constant, the elementary full von Mangoldt estimate in \cite{vfourTrace} gives $C_N\le16+6\sqrt N$. Put $L=1+\log(1/\delta)\ge1$ and $x=\sqrt L$. From \eqref{vfourD:Ndelta} and $\log2>1/2$, $N<8L$; hence $C_N<34x$ and $T<e^{34x}$. With $h=\max\{1,4T/\delta\}$, $\log h<36L$, so
\[R=\tfrac12\log\log(80h)<2+\tfrac12\log L\le\tfrac52x<e^{2x}.
\]
Also $\tfrac12\log_2(2h)<37L$. The explicit rank formula gives
\[r<3e^{36x}+37L+2\le42e^{36x}<e^{40x},
\]
using $L=x^2\le e^{36x}$ and $e^4>42$. This proves \eqref{vfourD:countbound}. The constants depend only on the displayed actual input estimates, not on a Gram condition number or an inverse computation.
\end{proof}
\begin{corollary}\label{vfourD:schatten}
The negative part $(L_c)_-=(-L_c)_+$ is compact and belongs to $\mathfrak S_p$ for every $p>0$, in particular to $\mathfrak S_1$.
\end{corollary}
\begin{proof}
The count is finite at every positive threshold, and $L_c$ is lower bounded. Its negative part is therefore the norm limit of finite-rank spectral truncations. If $\mu_n$ are its nonzero eigenvalues, the layer-cake formula for $p>0$ gives
\[\sum_n\mu_n^p=p\int_0^{\norm{(L_c)_-}}\delta^{p-1}\#\{n:\mu_n>\delta\}\,d\delta.
\]
The strict count under the integral is bounded by the closed count \eqref{vfourD:countbound}. For any fixed $p>0$, $e^{40\sqrt{1+\log(1/\delta)}}\le C_p\delta^{-p/2}$ near zero, and $\delta^{p/2-1}$ is integrable. For $0<p<1$, $\mathfrak S_p$ denotes the usual Schatten quasi-class. Trace class does not imply finitely many or no negative eigenvalues.
\end{proof}

